\section{Method}
\label{sec:methods}

\subsection{Q-transform maps}
\label{sec:methods:qscan}

Each trigger is analysed on a normalised-energy Q-transform map
\citep{Chatterji2004}. The transform tiles the time--frequency plane with
constant-$Q$ bisquare windows and reports, for each tile,
%
\begin{equation}
  Z = \frac{|X|^2}{\langle |X|^2 \rangle},
  \label{eq:normenergy}
\end{equation}
%
the ratio of tile energy to the mean over tiles at the same frequency. For
stationary Gaussian noise $Z$ is exponentially distributed with unit mean, so
$\sqrt{2Z}$ is the single-tile signal-to-noise ratio. The implementation is
\code{gwpy}'s \code{q\_transform} \citep{Macleod2021GWpy}, which is the same
algorithm used by the LIGO detector-characterisation tools; the standard
\code{gwdetchar} wrapper is not used because its dependency chain has no
Windows build, and nothing beyond the transform itself is needed.

\SI{64}{\second} of strain is fetched around each trigger and only the central
\SI{16}{\second} is analysed. The padding is not decorative: the whitening
filter and the constant-$Q$ filters both ring, and the resulting transients
are exactly the kind of structure a ridge tracker would mistake for an arch.
The configuration enforces that the padding exceeds the analysed window.

Two parameters were set by validation rather than by convention, and both
mattered.

\paragraph{The $Q$ range is capped at 16.} \code{gwpy} selects the $Q$ plane of
maximum energy. Left unconstrained it selects $Q \approx 45$ for a scattering
train, because a high-$Q$ plane integrates the quasi-periodic structure
coherently and reports enormous energy. But a tile of quality factor $Q$ at
frequency $f$ lasts $Q/2\pi f$ seconds --- \SI{0.16}{\second} at $Q = 45$ and
\SI{45}{\hertz} --- which is comparable to the repetition period of fast
scattering. The arches are then unresolved. On a simulated fast event with an
injected arch period of \SI{0.250}{\second}, the uncapped transform returned
\SI{4.99}{\second}: wrong by a factor of twenty. With the cap it returns
\SI{0.250}{\second}.

\paragraph{The ridge threshold is $Z = 20$, not $Z = 5.5$.} A threshold has to
sit above the largest energy a \emph{noise-only} map produces, not merely
above its mean. A map of this size holds of order $10^4$ independent tiles,
whose expected maximum is $\ln(10^4) \approx 9$ before the multi-$Q$ search is
taken into account and $\approx 16$ after. At $Z = 5.5$ --- a natural-looking
choice, corresponding to single-tile SNR $3.3$ --- a pure-noise map produced
hundreds of supra-threshold bins and the measured maximum arch frequency
saturated at the top of the analysis band. $Z = 20$ (single-tile SNR $6.3$)
clears the noise; a connected-component filter removes the residue.

\subsection{Ridge extraction}
\label{sec:methods:ridge}

From the thresholded map, three tracks are formed per time bin: the frequency
of the loudest tile, the highest supra-threshold frequency, and the lowest.
The second of these --- the \emph{arch envelope} --- is the quantity the
auxiliary predictor predicts, and is what the morphology features describe.

Connected supra-threshold regions smaller than $40$ output pixels are
discarded. Because the output grid is interpolated from far coarser $Q$ tiles,
a single hot tile already spans tens of pixels, so this is a speck filter
rather than a significance cut; the significance is carried by the threshold.

\subsection{Morphology features}
\label{sec:methods:features}

Each map is reduced to the twenty features listed in \cref{app:features}.
They fall into five groups: extent (peak, maximum and minimum frequency,
bandwidth, duration, SNR); arch structure (count, repetition period, period
scatter); harmonic structure (count and spacing at the apex); shape (RMS slope
and curvature of the envelope, ratio of rising to falling time); and energy
distribution (fraction below \SI{40}{\hertz}, fraction above
\SI{100}{\hertz}, spectral centroid, band concentration, time symmetry).

They are computed on the normalised-energy map, so they are insensitive to the
absolute calibration and directly comparable between detectors and between
runs. Where a feature is undefined --- \code{arch\_period} when fewer than two
arches are found --- it is returned as \code{NaN} rather than as a fabricated
value, and the classifier treats the missingness as informative.

Two of them deserve comment.

\paragraph{Arch period, from the autocorrelation.} The obvious estimator is
the median spacing between detected envelope maxima. It fails in exactly the
case that matters: when the topmost harmonic dips below threshold the envelope
drops to the harmonic below, creating a spurious local maximum that splits one
arch into two. On a simulated slow event with an injected period of
\SI{3.33}{\second}, peak spacing returned \SI{2.80}{\second} --- \SI{16}{\percent}
low. The autocorrelation of the envelope is insensitive to this, because the
periodicity is present in every harmonic simultaneously; it returns
\SI{3.36}{\second}. The peak count is still reported separately, as
\code{n\_arches}.

\paragraph{Maximum arch frequency is biased high, and stays that way.} A $Q$
tile has bandwidth $f/Q$, so the supra-threshold region extends above the true
apex of the arch by roughly half of that. Measured against injections, the
bias is about \SI{15}{\percent}. It is not tuned away: it is measured,
asserted in the test suite, and quoted with the result. Since it applies
identically to both runs it cancels in the O3-versus-O4a comparison, which is
where the feature is actually used.

\subsection{From triggers to events, and from events to blocks}
\label{sec:methods:sampling}

Two properties of the real data, neither visible until the catalogue and the
strain server were in hand, determine how the sample is built. Both are
recorded here because each one changes a number that would otherwise be
reported.

\paragraph{The catalogue lists triggers, not events.} Gravity Spy publishes
Omicron triggers. After the cuts of \cref{sec:data:gravityspy},
\SI{12.3}{\percent} of rows share an exact $(\mathrm{event\ time},
\mathrm{detector})$ pair -- the same trigger under several Gravity Spy subject
identifiers, one per Q-scan duration view -- and about half of all triggers
fall within \SI{8}{\second} of another one, which is to say inside each other's
analysed Q-scan window. Those triggers therefore produce \emph{identical}
feature vectors.

Left alone, every morphology distribution would be weighted by how many
Omicron triggers a given arch train happened to generate: loud, long trains
counted dozens of times, quiet ones once. That is a selection effect acting
directly on the axis this study measures. Worse, it is class-dependent.
Clustering at \SI{16}{\second} -- the full width of the analysed window, so
that no two retained events can share a Q-scan -- yields \num{2.44} triggers
per event for scattered light but only \num{1.35} for the control classes.
Uncorrected, scattering would have been over-represented by some
\SI{80}{\percent} relative to the controls, biasing both the classifier's
training set and every rate comparison. The same window is applied to the O4a
search output, so the two runs count the same kind of thing.

The unclustered trigger list is retained as well: the rate correlation of
\cref{sec:methods:ratecorr} legitimately uses \emph{trigger} rate, and each
result states which of the two it quotes.

\paragraph{The unit of cost is the file, not the event.} GWOSC serves open
strain in HDF5 files of \SI{4096}{\second}, and a request for any
sub-interval downloads the whole containing file: two to four minutes on a
domestic connection, essentially independent of how much was asked for. A
\SI{64}{\second} Q-scan window therefore costs one full download, and a
time-stratified sample of \num{4800} individual events touched \num{3484}
distinct files -- about \num{1.4} events per download, or three days of wall
clock.

The analysis is restructured around that. Windows are cut from whole blocks
held in a small in-memory cache, and the sample is drawn as \emph{blocks
first, events within blocks second}: \num{100} blocks per detector per epoch,
stratified across each epoch's span, with up to \num{12} events taken at
random from each. That gives \num{2256} scattered-light events from \num{400}
blocks, and the measured yield in a production run is about ten events per
download. The control sample is drawn from the \emph{same} blocks, which costs
no extra downloads and removes a confound at the same time: signal and
background then come from identical stretches of data, so the classifier
cannot separate the two classes by learning what a noisy night looks like.

The price is real and is carried into the interpretation. Events drawn from
one block share ground-motion conditions and are therefore correlated, so the
effective sample size is smaller than the row count: it is the \emph{block}
count, not the event count, that carries the statistical weight. Events are
taken at random within a block rather than by loudness, which would bias the
sample towards the tail.

\subsection{Rate-level auxiliary correlation}
\label{sec:methods:ratecorr}

Triggers are binned into GPS hours and divided by the analysis-ready livetime
in each hour. Hours with zero livetime and hours with zero triggers are
different things: the former must be dropped, the latter kept as genuine
zeros. Conflating them is a reliable way to manufacture a spurious correlation
with anything that tracks duty cycle, ground motion included.

Each hourly rate series is rank-correlated (Spearman $\rho$) against the
hourly median band-limited RMS of every configured seismometer channel. The
median within an hour is used rather than the mean because the trend channels
contain occasional large excursions --- earthquakes, calibration activity ---
whose effect on an hourly mean is out of proportion to how much of the hour
they occupy.

Fetching those trends has its own constraints, measured against the live
server rather than assumed. The GWOSC NDS host serves a four-hour request but
refuses eight; nine channels requested together cost \SI{11.3}{\second}
against \SI{40.8}{\second} one at a time, so requests are batched; and
requests are issued only for chunks that overlap analysis-ready segments,
since a chunk with no livetime contributes nothing to a rate correlation.
Availability looked poor when probed at fixed calendar offsets and turned out
to be complete -- 18 of 18 -- when probed uniformly in livetime, the earlier
failures having been non-observing time.

Confidence intervals come from a moving-block bootstrap
\citep{Kunsch1989} with a six-hour block. This matters: the microseism varies
on a timescale of days, so hourly rates are strongly autocorrelated and the
naive standard error on $\rho$ is far too small. Three further checks guard
against the standard failure modes:

\begin{itemize}
  \item a \textbf{quantile table}, giving the mean rate in each quartile of
        the driving channel and the top-to-bottom ratio, so that the result is
        stated in units a commissioner can act on rather than as a correlation
        coefficient;
  \item a \textbf{lag scan} over $\pm\SI{12}{\hour}$. Scattering responds to
        ground motion essentially instantaneously, so the correlation must
        peak at zero lag. A peak at large lag would indicate that both series
        are being driven by a slow common trend --- diurnal or weather ---
        rather than by the mechanism under test;
  \item a \textbf{partial correlation}, testing whether an apparent
        correlation with the \SIrange{1}{6}{\hertz} anthropogenic band
        survives once the microseism is controlled for.
\end{itemize}

\subsection{Trigger-level scattering predictor}
\label{sec:methods:predictor}

For \Lone, where a full-rate motion witness is public, the predictor of
\cref{eq:fn} is tested trigger by trigger.

The witness is band-passed over \SIrange{0.03}{10}{\hertz} --- the lower
corner removes the DC drift that differentiation would otherwise amplify into
a ramp --- and differentiated to give velocity. Harmonics $n = 1 \ldots 8$ of
\cref{eq:fn} are then compared with the measured arch envelope. A ridge point
at frequency $f$ counts as \emph{explained} if some harmonic lies within
$\max(\SI{1.5}{\hertz},\, 0.15 f)$ of it, and the headline statistic is the
fraction of supra-threshold time bins that are explained.

\paragraph{The one free parameter.} GWOSC publishes auxiliary channels in raw
counts with no calibration to metres. Rather than invent one, a single scale
factor $\alpha$ (metres per count) is fitted per detector and epoch by
maximising the mean explained fraction over a training set of triggers, using
a logarithmic grid over eight decades followed by a golden-section refinement.
It is then \emph{frozen} and applied to a disjoint validation set, on which the
quoted agreement is measured. The width of the range over which the objective
stays within \SI{1}{\percent} of its maximum is reported alongside: a narrow
range means the fit is well constrained, a wide one is a warning that the
witness does not in fact track the scattering surface.

\Cref{sec:validation} shows this procedure recovering a known calibration from
a synthetic witness to \SI{7}{\percent}.

\subsection{The strain-only identifier}
\label{sec:methods:classifier}

To carry the O3 definition of the class into O4a, a gradient-boosted tree
classifier \citep{Friedman2001,Pedregosa2011} is trained on the O3 features to
separate \emph{Scattered\_Light} from the control classes. Missing features are
median-imputed with a missingness indicator, so that ``no arch period could be
measured'' is available to the model as information rather than being silently
filled in.

Performance is measured by stratified five-fold cross-validation, and the
operating point is fixed at a false-positive rate of $0.01$ against the
control sample, then frozen.

A third constraint is arithmetic and binds harder than either of the caveats
below. An operating point can only be placed as finely as the background
sample resolves: at a false-positive rate $p$ with $n$ control events per fold,
the threshold is fixed by $p\,n$ expected false positives. Trained on a trial
sample of 66 controls, the cross-validated AUC was a stable
$0.919 \pm 0.033$ while the efficiency at $p = 0.01$ came out
$0.368 \pm 0.334$, with individual folds returning $0.000$ (at an infinite
threshold), $0.160$, $0.680$, $0.760$ and $0.240$ --- because $p\,n = 0.13$,
so the threshold was decided by whether zero or one event happened to land
above it. The code now computes $p\,n$, warns below five, and records the fact
in the saved model; the production control sample gives $p\,n \approx 7$ and
clears it. The AUC is threshold-free and unaffected throughout.

Two further caveats travel with those numbers and are repeated wherever they
are quoted. The efficiency and false-positive rate are measured against the
\emph{Gravity Spy control classes}, so they say how well the model separates
scattered light from other known O3 glitch classes; they do not and cannot
measure performance against whatever unlabelled O4a contains. And the folds
are not grouped in time, so triggers minutes apart can land in different
folds, making the quoted AUC mildly optimistic. The sample is drawn stratified
across the whole run, which makes such neighbours rare but does not eliminate
them.

\subsection{The O4a search}
\label{sec:methods:search}

With no O4a trigger catalogue, triggers are generated from strain. The
statistic is the \SI{0.25}{\second} RMS of whitened, band-passed
(\SIrange{10}{120}{\hertz}) strain, converted to a robust $z$-score using the
median and median-absolute-deviation of the surrounding stretch. Robust
statistics are essential here: during a scattering-heavy hour a large fraction
of the data is glitching, and a mean/standard-deviation normalisation would be
dragged upward by the very signal it is meant to detect. Candidates are
clustered with a \SI{4}{\second} window, since an arch train is one physical
event lasting seconds, and candidates within the whitening-filter transient at
either end are discarded.

This is a deliberately loose stand-in for Omicron: its job is to lose no
scattering, not to be pure. Purity is the classifier's job.

\paragraph{Livetime, not calendar time.} A full O4a search is out of reach on
the hardware available. The search draws $1200$ windows of \SI{512}{\second}
per detector, sampled uniformly in livetime --- about \SI{7.1}{\day} each,
roughly \SI{4}{\percent} of the run. All rates are therefore quoted per hour of
\emph{analysed} livetime, and the coverage fraction is recorded with every
result.

\paragraph{Like-for-like comparison.} The identical search is also run on O3b.
This is what makes the comparison interpretable: with the same trigger
generator and the same classifier on both sides, a difference in rate or
morphology cannot be attributed to the change from the Gravity Spy pipeline to
this one. The comparison against the Gravity Spy O3 sample is also reported,
since it has better labels; where the two agree the result is robust, and
where they disagree the difference is a selection effect and is reported as
such.

\subsection{Population statistics}
\label{sec:methods:stats}

Rates are Poisson counts over a known livetime, so intervals are exact
(Garwood) Poisson intervals rather than $\sqrt{N}$.

Feature distributions are compared with the two-sample Kolmogorov--Smirnov
test for the $p$-value and with Cliff's $\delta$ \citep{Cliff1993} for the
effect size. With tens of thousands of triggers per sample essentially every
KS test rejects, so the $p$-value carries little information; the effect size
is what carries the physics and is quoted first. Because many features are
tested at once, $p$-values are corrected by the Benjamini--Hochberg procedure
\citep{BenjaminiHochberg1995} and the corrected values are the ones reported.

The headline slow-fraction change is a difference of two binomial proportions,
reported with a Newcombe score-based interval \citep{Newcombe1998}, which
behaves correctly as a proportion approaches $0$ or $1$ where the Wald
interval does not.
